%% kapa: exam=1 round=28 year=2017 period=2 form=A subject=appraisal_law subject_ko=감정평가관계법규 seq=38 number=38
%% source: Q-Net 감정평가사 기출문제 (cst00309) artlSeq=5199603 file=extracted/1008/5199603/2017년 제28회 감정평가사 1차시험 문제지/2교시 A형.pdf
%% conversion: kordoc 4.24.51 -> tools/gen1.py (auto); review=proofread-sonnet; flags=none
\begin{question}{38}
도시 및 주거환경정비법령상 관리처분계획에 관한 설명으로 옳은 것은?
\begin{choices}
  \item 주택재건축사업에서 주택분양에 관한 권리를 포기하는 토지등소유자에 대한 임대 주택의 공급에 따라 관리처분계획을 변경하는 때에는 시장ㆍ군수에게 신고하여야 한다.
  \item 주택재개발사업에서 관리처분계획은 주택단지의 경우 1개의 건축물의 대지는 1필지의 토지가 되도록 정하여야 한다.
  \item 주택재건축사업의 관리처분계획에서 분양대상자별 분양예정인 건축물의 추산액을 평가할 때에는 시장ㆍ군수가 선정ㆍ계약한 2인 이상의 감정평가업자가 평가한 금액을 산술평균하여 산정한다.
  \item 주택재개발사업에서 지방자치단체인 토지등소유자에게는 하나 이상의 주택 또는 토지를 소유한 경우라도 1주택을 공급하도록 관리처분계획을 정한다.
  \item 주거환경관리사업의 사업시행자는 관리처분계획을 수립하여 시장ㆍ군수의 인가를 받아야 한다.
\end{choices}
\end{question}
\answer{1}
